\section*{अध्याय \ref{ch:g12:mountain-building} --- पर्वतों का निर्माण}

\begin{solution}{exo:g12:mountain-building:1}
(a) \omterm{def:g12:mountain-building:thrust}{प्रणोद}; (b) \omterm{def:g12:mountain-building:ophiolite}{ओफ़ियोलाइट}; (c) \omterm{def:g12:mountain-building:fold}{अभिनति}; (d) \omterm{def:g12:mountain-building:suture}{सीवन}; (e) \omterm{def:g12:mountain-building:fold}{अपनति};
(f) \omterm{def:g12:mountain-building:nappe}{नैप}।
\end{solution}

\begin{solution}{exo:g12:mountain-building:2}
(1) भुजाएँ \omterm{def:g12:mountain-building:fold}{वलन} के अक्ष से दूर नत हैं: \omterm{def:g12:mountain-building:fold}{अपनति}। (2) सबसे पुरानी
परतें। (3) नहीं: अपरदन कठोर और नरम परतों का अनुसरण करता है, \omterm{def:g12:mountain-building:fold}{वलन} की
आकृति का नहीं; \omterm{def:g12:mountain-building:fold}{अपनति} का क्रोड घिसकर घाटी बन सकता है।
\end{solution}

\begin{solution}{exo:g12:mountain-building:3}
(1) $e = 1 - \cos\ang{30} = 1 - 0.866 = 0.134$, लगभग \qty{13}{\percent}।
(2) $L_0 = 6.0/\cos\ang{30} \approx \qty{6.9}{km}$।
\end{solution}

\begin{solution}{exo:g12:mountain-building:4}
(1) \omterm{def:g11:ocean-ridges:normal-fault}{सामान्य भ्रंश}; भूपर्पटी खिंची थी। (2) हाँ: यह \ang{45} से कम नत
\omterm{def:g12:mountain-building:thrust}{उत्क्रम भ्रंश} है।
\end{solution}

\begin{solution}{exo:g12:mountain-building:5}
(1) समुद्री अवसाद, तकिया बेसाल्ट, चादरी \omterm{def:g11:igneous-rocks:intrusion}{डाइक}, \omterm{def:g11:ocean-ridges:gabbro}{गैब्रो}, \omterm{def:g11:ocean-ridges:peridotite}{पेरिडोटाइट}।
(2) लावा तकियों का रूप तभी लेता है जब वह पानी में ठंडा होता है। (3) \omterm{def:g12:mountain-building:ophiolite}{ओफ़ियोलाइट}
उस महासागर का एक टुकड़ा है जो दो महाद्वीपों के बीच था; जब वे टकराए,
तो उसकी तली की फाँकें डूबने के बजाय उनके किनारों पर धकेल दी गईं, इसलिए वे
वहाँ मिलती हैं जहाँ महासागर बंद हुआ, \omterm{def:g12:mountain-building:suture}{सीवन} के अनुदिश।
\end{solution}

\begin{solution}{exo:g12:mountain-building:6}
(1) $L_0 = 2.1 + 1.8 + 2.4 + 1.6 = \qty{7.9}{km}$;
$e = (7.9 - 6.0)/7.9 \approx 0.24$, \qty{24}{\percent}। (2) $\cos\theta =
6.0/7.9 \approx 0.759$, $\theta \approx \ang{41}$।
\end{solution}

\begin{solution}{exo:g12:mountain-building:7}
(1) $45\cos\ang{10} \approx \qty{44.3}{km}$। (2) $45\sin\ang{10} \approx
\qty{7.8}{km}$। (3) ऐसे \omterm{def:g12:mountain-building:fold}{वलन} $1 - \cos\ang{12} \approx 0.0219$ से छोटा करते हैं:
\omterm{def:g12:mountain-building:fold}{वलन} से पहले पट्टी $44.3/0.0219 \approx \qty{2000}{km}$ चौड़ी रही
होगी।
\end{solution}

\begin{solution}{exo:g12:mountain-building:8}
(1) अध्यारोपण: भ्रंश को अनदेखा करके पढ़ने पर ऊपरी चट्टान को
नई मान लिया जाएगा। (2) $500 - 75 = 425$ मिलियन वर्ष। (3) $\qty{40000}{m} /
\qty{800}{m} = 50$ गुना।
\end{solution}

\begin{solution}{exo:g12:mountain-building:9}
(1) $(30.23 - 19.38) \times 111.2 = 10.85 \times 111.2 \approx
\qty{1207}{km}$। (2) $1207/14 \approx \qty{86}{km}$ प्रति मिलियन वर्ष,
यानी प्रति वर्ष \qty{86}{mm}, लगभग \qty{8.6}{cm}। (3) प्रति शताब्दी \qty{9}{m}
का अर्थ है प्रति वर्ष \qty{9}{cm}: दोनों लगभग \qty{5}{\percent} के भीतर सहमत हैं।
\end{solution}

\begin{solution}{exo:g12:mountain-building:10}
(1) $H = 35/(1 - 0.40) \approx \qty{58}{km}$। (2) $\Delta H \approx
\qty{23.3}{km}$; वह $23.3 \times 550/3380 \approx
\qty{3.8}{km}$ अधिक ऊँची खड़ी होती है, और उसका अतिरिक्त \omterm{def:g12:continental-crust:root}{मूल} $23.3 - 3.8 \approx \qty{19.5}{km}$ गहरा है।
\end{solution}

\begin{solution}{exo:g12:mountain-building:11}
(1) चित्र $\tan\theta$ को 3 से गुणा कर देता है: $\tan\theta = \tan\ang{40}/3
\approx 0.280$, $\theta \approx \ang{15.6}$। (2) वास्तविक: $1 - \cos\ang{15.6}
\approx 0.037$, \qty{3.7}{\percent}; चित्र पर पढ़ा गया: $1 - \cos\ang{40}
\approx 0.234$, \qty{23}{\percent}। (3) खींचा गया ऊर्ध्वाधर पैमाना
नमनों को अधिक खड़ा और परतों को अधिक लंबा कर देता है, इसलिए वह \omterm{def:g12:mountain-building:shortening}{लघुकरण} को कई
गुना अधिक बड़ा दिखाता है।
\end{solution}

\begin{solution}{exo:g12:mountain-building:12}
(1) $L_0 = 10/(1 - 0.317) \approx \qty{14.6}{km}$। (2) $1 - \cos\ang{57.5}
\approx 1 - 0.537 = 0.463$, \qty{46}{\percent}। (3) गोल \omterm{def:g12:mountain-building:fold}{वलन} में
परत का केवल एक छोटा भाग \ang{57.5} पर नत होता है; कब्ज़ों के पास वह
लगभग समतल है, और समतल भाग छोटे नहीं होते। सीधी भुजाओं वाले \omterm{def:g12:mountain-building:fold}{वलन} में
पूरी परत \ang{57.5} पर नत होती है।
\end{solution}

\begin{solution}{exo:g12:mountain-building:13}
(1) $\sqrt{16.6^2 + 41.8^2} \approx \qty{45.0}{mm/yr}$। (2) $41.8 - 23.5 =
\qty{18.3}{mm}$ प्रति वर्ष, यानी ल्हासा के उत्तर में प्रति मिलियन वर्ष लगभग
\qty{18}{km} उत्तर--दक्षिण \omterm{def:g12:mountain-building:shortening}{लघुकरण}। (3) स्वयं ल्हासा उत्तर की ओर चला है, और
पट्टी के भूकंप हिमालय से बहुत उत्तर तक एशिया के आर-पार फैले हैं: वहाँ की
भूपर्पटी अब भी टूट रही है।
\end{solution}

\begin{solution}{exo:g12:mountain-building:14}
(1) उसका नमन लगभग शून्य है, इसलिए \omterm{def:g12:mountain-building:shortening}{लघुकरण} ही विस्थापन है,
\qty{45}{km}। (2) $45/(1 - \cos\ang{12}) = 45/0.0219 \approx \qty{2060}{km}$।
(3) बड़े अभिसरण मुख्यतः \omterm{def:g12:mountain-building:thrust}{प्रणोदों} और \omterm{def:g12:mountain-building:nappe}{नैपों} से, और कसे हुए \omterm{def:g12:mountain-building:fold}{वलनों} से,
अवशोषित होते हैं; मंद \omterm{def:g12:mountain-building:fold}{वलन} बहुत कम अवशोषित करते हैं।
\end{solution}

\begin{solution}{exo:g12:mountain-building:15}
(1) $\qty{8848.86}{m} = \qty{8848860}{mm}$, $\num{5e7}$ वर्षों में: लगभग
प्रति वर्ष \qty{0.18}{mm}। (2) $40/0.18 \approx 220$: उत्थान भारत की चाल से दो सौ
गुना से भी अधिक धीमा है। (3) चूना पत्थर समुद्र-तल से नीचे जमा था,
और अपरदन ने उसके ऊपर पड़ी चट्टानें हटा दी हैं: वह अपनी वर्तमान ऊँचाई से
अधिक उठाया गया है।
\end{solution}

\begin{solution}{pb:g12:mountain-building:1}
\textbf{1.} $2\pi \times 6371/360 \approx \qty{111.2}{km}$।

\textbf{2.} 50 मिलियन वर्ष पहले (\ang{0.68}~द) और 45 (\ang{3.30}~उ) के बीच:
$50 - 5 \times 0.68/3.98 \approx 49$ मिलियन वर्ष पहले।

\textbf{3.} $(19.38 - 0.68) \times 111.2 = 18.70 \times 111.2 \approx
\qty{2079}{km}$।

\textbf{4.} $2079/16 \approx \qty{130}{km}$ प्रति मिलियन वर्ष: प्रति वर्ष \qty{130}{mm},
लगभग \qty{13}{cm}।

\textbf{5.} 66 से 60: $7.28 \times 111.2/6 \approx 135$; 60 से 55: $4.86
\times 111.2/5 \approx 108$; 55 से 50: $6.56 \times 111.2/5 \approx 146$; 50
से 45: $3.98 \times 111.2/5 \approx 89$; 45 से 40: $2.29 \times 111.2/5
\approx 51$ (प्रति वर्ष \unit{mm})। सबसे तेज़ 55 से 50 मिलियन वर्ष पहले तक।

\textbf{6.} $(20.00 - 5.59) \times 111.2/40 = 14.41 \times 111.2/40 \approx
\qty{40}{mm}$ प्रति वर्ष।

\textbf{7.} $40/130 \approx 0.31$: एक-तिहाई से कम।

\textbf{8.} चरण 55--50 (146) से चरण 50--45 (89) तक: लगभग 50
मिलियन वर्ष पहले।

\textbf{9.} \omterm{def:g12:continental-crust:continental-crust}{महाद्वीपीय भूपर्पटी} डूबने के लिए बहुत हल्की है: जब वह
खाई तक पहुँची, तो \omterm{def:g11:subduction-zones:subduction}{क्षेपण} भारत को खुलकर साथ नहीं खींच सका, और दूसरे
महाद्वीप ने उसे रोक लिया।

\textbf{10.} $(20.00 + 0.68) \times 111.2 = 20.68 \times 111.2 \approx
\qty{2300}{km}$।

\textbf{11.} $(29.65 - 21.15) \times 111.2 = 8.50 \times 111.2 \approx
\qty{945}{km}$।

\textbf{12.} $2300 - 945 \approx \qty{1354}{km}$।

\textbf{13.} आज $(29.65 - 20.00) \times 111.2 \approx \qty{1073}{km}$;
50 मिलियन वर्ष पहले $(21.15 + 0.68) \times 111.2 = 21.83 \times 111.2
\approx \qty{2427}{km}$ (और $2427 - 1073 = 1354$)।

\textbf{14.} $e = 1354/2427 \approx 0.56$, \qty{56}{\percent}।

\textbf{15.} \omterm{def:g12:mountain-building:fold}{वलन}, और \omterm{def:g12:mountain-building:thrust}{उत्क्रम भ्रंश}: \omterm{def:g12:mountain-building:nappe}{नैपों} को ढोने वाले \omterm{def:g12:mountain-building:thrust}{प्रणोद}।

\textbf{16.} $H = 35 \times 2427/1073 \approx \qty{79}{km}$।

\textbf{17.} $\Delta H = 79 - 35 = \qty{44}{km}$; वह
$44 \times 550/3380 \approx \qty{7.2}{km}$ अधिक ऊँची खड़ी होगी।

\textbf{18.} $\Delta H = 4.6 \times 3380/550 \approx \qty{28}{km}$, यानी
$35 + 28 = \qty{63}{km}$ मोटी भूपर्पटी, मापे गए \qty{70}{km} के
निकट।

\textbf{19.} भूपर्पटी अपनी अनुप्रस्थ काट बनाए नहीं रखती: एशिया के कुछ भाग
पूर्व की ओर भींचकर बाहर निकाले जाते हैं, और अपरदन पर्वतों के शीर्ष से
चट्टान हटाता है।

\textbf{20.} $2079 + 2300 \approx \qty{4380}{km}$: क्रिटेशस के अंत से
भारत लगभग \textbf{\qty{4400}{km} उत्तर की ओर} यात्रा कर चुका है।
\end{solution}
